\documentclass[11pt]{article}

\usepackage[margin=1.2in]{geometry}
\usepackage{booktabs}
\usepackage{graphicx}
\usepackage{hyperref}
\usepackage{amsmath}
\usepackage{array}
\setlength{\parskip}{6pt}
\setlength{\parindent}{0pt}

\hypersetup{colorlinks=true, linkcolor=black, urlcolor=blue, citecolor=black}

\title{Lens Patent Dataset Overview\\
\large What we loaded, how we cleaned it, and what's in it}
\author{Mantis IP Visualization}
\date{May 2026}

\begin{document}
\maketitle
\tableofcontents
\newpage

% ─────────────────────────────────────────────────────────────────────────────
\section{What is the Lens Dataset?}
% ─────────────────────────────────────────────────────────────────────────────

Lens.org is a free, open-access patent search platform maintained by Cambia, a non-profit
organization focused on democratizing access to innovation. It aggregates patent records from
major offices worldwide—including the USPTO, EPO, WIPO, JPO, and CNIPA—covering over
150 million patents and patent applications. Each record includes structured bibliographic
metadata, legal status, family linkages, classification codes, and, where available, full text
fields such as the abstract, description, and claims. The Lens API allows bulk export of
structured JSON records, making it well-suited for downstream machine learning tasks.

% ─────────────────────────────────────────────────────────────────────────────
\section{What's in the raw dataset}
% ─────────────────────────────────────────────────────────────────────────────

The raw export contains \textbf{8,357 records} drawn from the Lens database via a targeted
query on IP relevant to the Mantis visualization project. Records span 43 distinct
jurisdictions and 20 languages, and cover publications from the late 19th century through
2024, though the overwhelming majority were published after 2017. The dataset is intentionally
broad and largely uncleaned at this stage.

\subsection{Jurisdictions and languages}

The raw pull spans 43 jurisdictions and 20 languages. Non-US records make up the majority,
with China, Japan, and the EPO being the largest non-US sources. The multilingual character
is a direct consequence of the broad query scope and is resolved during cleaning
(Section~\ref{sec:cleaning}).

\vspace{1em}
\begin{figure}[h]
  \centering
  \includegraphics[width=0.92\textwidth]{figures/fig1_jurisdiction.png}
  \caption{Jurisdiction distribution across the raw dataset (top 12 offices shown, n\,=\,8,357).}
  \label{fig:jurisdiction}
\end{figure}

\vspace{1.5em}
\begin{figure}[h]
  \centering
  \includegraphics[width=0.92\textwidth]{figures/fig2_language.png}
  \caption{Language distribution across the raw dataset.}
  \label{fig:language}
\end{figure}
\vspace{1em}

\subsection{Field completeness}

Text field coverage is uneven. The abstract is the most reliably present field, while
description and claims are missing from over half of records—typical for non-US or older
filings. The sequence listing field is not semantically relevant and was dropped during
cleaning.

\vspace{1em}
\begin{figure}[h]
  \centering
  \includegraphics[width=0.92\textwidth]{figures/fig4_field_completeness.png}
  \caption{Percentage of raw records with each field populated.}
  \label{fig:completeness}
\end{figure}
\vspace{1em}

\subsection{Publication types and legal status}

The dataset is a mix of applications and granted patents, with a long tail of less common
types. Most patents are either active or pending, reflecting the recency of the dataset.

\vspace{1em}
\begin{figure}[h]
  \centering
  \includegraphics[width=0.92\textwidth]{figures/fig5_publication_type.png}
  \caption{Publication type distribution across the raw dataset.}
  \label{fig:pubtype}
\end{figure}

\vspace{1.5em}
\begin{figure}[h]
  \centering
  \includegraphics[width=0.92\textwidth]{figures/fig6_patent_status.png}
  \caption{Legal status distribution across the raw dataset.}
  \label{fig:status}
\end{figure}
\vspace{1em}

\subsection{Publication year}

The dataset is heavily skewed toward recent filings, with the bulk of records from 2020
onward. There is a thin tail of historical records extending back to the 19th century.

\vspace{1em}
\begin{figure}[h]
  \centering
  \includegraphics[width=0.92\textwidth]{figures/fig7_publication_year.png}
  \caption{Publication year distribution (raw dataset, 2000--2024).}
  \label{fig:year}
\end{figure}
\vspace{1em}

\subsection{Technology coverage (IPC sections)}

All eight WIPO top-level sections are represented. The heaviest concentration is in
Physics (G) and Electricity (H), consistent with the technology domains most relevant to Mantis.

\vspace{1em}
\begin{figure}[h]
  \centering
  \includegraphics[width=0.92\textwidth]{figures/fig8_ipc_sections.png}
  \caption{Number of unique patents per IPC top-level section. A patent may appear in
  multiple sections.}
  \label{fig:ipc}
\end{figure}
\vspace{1em}

% ─────────────────────────────────────────────────────────────────────────────
\section{How we cleaned the data}
\label{sec:cleaning}
% ─────────────────────────────────────────────────────────────────────────────

The raw export was too broad and inconsistent for direct use in Mantis. Three sequential
filters were applied, followed by a structural flattening of the nested JSON schema.

\subsection{Filtering steps}

\begin{enumerate}
  \item \textbf{US jurisdiction only.} Records from non-US patent offices were removed.
    This reduced the dataset from 8,357 to 2,818 records and incidentally resolved the
    multilingual issue: after this filter, 99.9\% of remaining records are in English, with
    three records carrying no language tag (presumed English given US jurisdiction).

  \item \textbf{Abstract required.} Records with no abstract field were dropped. The
    abstract is the primary semantic signal used by Mantis; records without one are not
    useful. After the jurisdiction filter, only 39 of the 2,818 US records lacked an
    abstract, leaving 2,779.

  \item \textbf{Deduplication.} Records were deduplicated on \texttt{lens\_id}. No
    duplicates were found in the US-filtered subset, so the count remained at 2,779.
\end{enumerate}

\subsection{Schema flattening}

The raw Lens records use a deeply nested JSON structure (e.g., inventor names nested inside
\texttt{biblio\,>\,parties\,>\,inventors\,>\,extracted\_name\,>\,value}). All records were
flattened to a single-level schema with 19 columns. Fields that were either redundant,
internal database identifiers, or not relevant for patent identification or semantic analysis
were dropped during this step. Any column missing from a given record was filled with the
string \texttt{"null"} to ensure a consistent schema across all 2,779 samples, as required
by Mantis.

% ─────────────────────────────────────────────────────────────────────────────
\section{Final dataset at a glance}
\label{sec:final}
% ─────────────────────────────────────────────────────────────────────────────

\begin{table}[h]
  \centering
  \begin{tabular}{ll}
    \toprule
    \textbf{Property} & \textbf{Value} \\
    \midrule
    Total records           & 2,779 \\
    Jurisdiction            & US only \\
    Language                & English (99.9\%), unknown (0.1\%) \\
    Fields per record       & 19 \\
    Records with abstract   & 2,779 (100\%) \\
    Records with description & 1,196 (43.0\%) \\
    Abstract: mean length   & 105 words \\
    Abstract: median length & 107 words \\
    Abstract: 90th pctile  & 149 words \\
    \bottomrule
  \end{tabular}
  \caption{Summary statistics for the cleaned dataset.}
  \label{tab:summary}
\end{table}

% ─────────────────────────────────────────────────────────────────────────────
\section{Schema reference}
\label{sec:schema}
% ─────────────────────────────────────────────────────────────────────────────

Table~\ref{tab:schema} describes each of the 19 columns in the cleaned dataset.

\begin{table}[h]
  \centering
  \small
  \begin{tabular}{p{4.8cm} p{7.5cm}}
    \toprule
    \textbf{Field} & \textbf{Description} \\
    \midrule
    \texttt{lens\_id}              & Globally unique Lens.org identifier for the patent record. \\
    \texttt{jurisdiction}          & Patent office jurisdiction (always \texttt{US} in this dataset). \\
    \texttt{doc\_number}           & USPTO document number. \\
    \texttt{kind}                  & Document kind code indicating publication stage (e.g., \texttt{A}, \texttt{B2}). \\
    \texttt{date\_published}       & Date the document was published (YYYY-MM-DD). \\
    \texttt{lang}                  & Language of the document text. \\
    \texttt{publication\_type}     & Lens-assigned publication category (e.g., \texttt{PATENT\_APPLICATION}, \texttt{GRANTED\_PATENT}). \\
    \texttt{patent\_status}        & Current legal status (e.g., \texttt{ACTIVE}, \texttt{PENDING}, \texttt{EXPIRED}). \\
    \texttt{title}                 & Invention title(s) as a list; may include translations. \\
    \texttt{application\_date}     & Date the application was filed with the USPTO. \\
    \texttt{priority\_date}        & Date of the earliest priority claim. \\
    \texttt{applicants}            & List of applicant (assignee) names. \\
    \texttt{inventors}             & List of inventor names. \\
    \texttt{ipc\_classifications}  & List of IPC classification symbols assigned to the patent. \\
    \texttt{cited\_by\_lens\_ids}  & Lens IDs of patents that cite this record (forward citations). \\
    \texttt{simple\_family\_lens\_ids}   & Lens IDs of members of the simple patent family. \\
    \texttt{extended\_family\_lens\_ids} & Lens IDs of extended patent family members. \\
    \texttt{abstract}              & Full abstract text as filed. \\
    \texttt{description}           & Full description text, where available. \\
    \bottomrule
  \end{tabular}
  \caption{Schema of the cleaned dataset (19 fields). Fields with no value for a given
  record are stored as the string \texttt{"null"}.}
  \label{tab:schema}
\end{table}

\end{document}
